\color{unidarkgreen}{\subsection{Фиксация положения выключателей (ФПВ)}}
\color{black}

\begin{enumerate}[label=\arabic{section}.\arabic{subsection}.\arabic*,  align=left, labelsep=4pt, leftmargin=0pt, itemindent=57pt]

\item
Логика ФПВ В1(ЛВ), ФПВ В2(ОВ), ФПВ В3 аналогичны, рассмотрим на примере ФПВ В1(ЛВ).

\uchyph=1
\item
Ввод в работу ФПВ В1(ЛВ) и осуществляется с помощью уставок \textbf{«ФПВ В1(ЛВ)»}. Оперативный
ввод/вывод ФПВ производится с помощью БУ (блоков управления функциями) \textit{«Вывод ФПВ В1(ЛВ)»}.
Предусматривается возможность формирования положения выключателей с контролем
разъединителей, которая задается уставками \textbf{«Контроль Р В1(ЛВ)»}. С помощью уставок \textbf{«Схема Р
В1(ЛВ)»} имеется возможность выбрать схему выключателя, два или три разъединителя. С помощью
уставки \textbf{«Тип В1(ЛВ)»} имеется возможность выбрать тип выключателя, линейный или обходной. Так
же от блок-контактов В1(ЛВ) формируется  внутренний пуск ЗНФ, который пофазно контролирует
расхождение контактов полюсов выключателей, возникающее при трехфазном включении или
отключении выключателя.
\uchyph=0

\begin{figure}[H]
\centering
\includegraphics[width=0.65\textwidth,height=0.65\textheight,keepaspectratio]{\fbpath/761.3471 ФПВ/img/fpv.pdf}
\captionof{figure}{Функциональная схема ФПВ В1}\label{pict_fpv:FPV}
\end{figure}

\small
\begin{longtable}{|>{\centering\arraybackslash}m{5.3cm}|>{\centering\arraybackslash}m{4.2cm}|>{\centering\arraybackslash}m{1cm}|>{\centering\arraybackslash}m{3.3cm}|>{\centering\arraybackslash}m{1.8cm}|}
% Первая страница: полный заголовок
\caption{Перечень уставочных параметров ФПВ\hfill\vspace{-0.5\baselineskip}}\label{fpv:tbl1}\\
\hline
\rowcolor{gray!30}
Наименование уставки & Диапазон & Шаг & Значение по умолчанию & Единица измерения \\
\hline
\endfirsthead

% Промежуточные страницы: «Продолжение таблицы»
\caption*{\hspace{3pt}\emph{Продолжение таблицы \ref{fpv:tbl1}}\hfill\ }\vspace{-0.5\baselineskip} \\
\hline
\rowcolor{gray!30}
Наименование уставки & Диапазон & Шаг & Значение по умолчанию & Единица измерения \\
\hline
\endhead

%===m>ФПВ|-
%===m
\end{longtable}
\normalsize

\end{enumerate}